\chapter{规范群与标准模型表示结构}
\label{chap:lie-standard-model}

前面的群都可以先当作全局对称：同一个群元在所有时空点共同作用。若允许变换参数随位置改变，普通导数会额外作用到参数上，场方程便不再协变。引入取值于 Lie 代数的联络并用协变导数补偿这个额外项，局域对称才真正成立；这就是规范场出现的表示论原因。

一代费米子和 Higgs 场各自承载 \(SU(3)_c\times SU(2)_L\times U(1)_Y\) 的一个表示。协变导数说明三个因子的生成元怎样分别作用，\(Q=T_3+Y\) 从权中读出电荷，张量积里的平凡表示则决定一个耦合能否写进拉格朗日量。最后还要把所有左手 Weyl 场放在一起检查反常；单个多重态合法，并不保证整套量子理论一致。

\section{局域变换为什么需要规范联络}

先从一个一般矩阵群 $G$ 的表示开始。物质场按
\[
\psi(x)\longmapsto\psi'(x)=U(x)\psi(x)
\]
变换，其中 $U(x)$ 可以随时空点改变。普通导数给出
\[
\partial_\mu\psi'
=U\,\partial_\mu\psi+(\partial_\mu U)\psi.
\]
第二项含有变换参数的导数，所以 $\partial_\mu\psi$ 不再按与
$\psi$ 相同的方式变换。引入 Lie 代数值场
$A_\mu=A_\mu^aT^a$，定义
\begin{equation}
D_\mu=\partial_\mu-\ii gA_\mu .
\label{eq:gauge-covariant-derivative-general}
\end{equation}
要求 $D_\mu'\psi'=U D_\mu\psi$，逐项比较得到
\begin{equation}
A_\mu'
=UA_\mu U^{-1}
-\frac{\ii}{g}(\partial_\mu U)U^{-1}.
\label{eq:gauge-connection-transformation}
\end{equation}
非齐次的第二项正好抵消 $\partial_\mu U$。规范势因此不按普通
张量变换；它的任务是比较相邻时空点处属于不同局域基的场。

两次协变微分的交换子不再含 $\psi$ 的导数：
\[
[D_\mu,D_\nu]\psi=-\ii gF_{\mu\nu}\psi,
\]
其中
\begin{equation}
F_{\mu\nu}
=\partial_\mu A_\nu-\partial_\nu A_\mu
-\ii g[A_\mu,A_\nu].
\label{eq:nonabelian-field-strength}
\end{equation}
利用\eqnrefs{eq:gauge-connection-transformation}直接代入可得
$F_{\mu\nu}'=UF_{\mu\nu}U^{-1}$。因此
$\Tr(F_{\mu\nu}F^{\mu\nu})$ 在规范变换下不变。对 $U(1)$，
所有生成元对易，交换子项消失，\eqnrefs{eq:nonabelian-field-strength}
退化为熟悉的电磁场强；对非 Abel 群，规范场本身带群指标，
交换子项产生规范玻色子的自相互作用。

\section{标准模型的表示结构}

标准模型的局域内部对称群在 Lie 代数层面为（式 \eqnrefs{eq:standard-model-group}）
\begin{equation}
G_{\mathrm{SM}}=SU(3)_c\times SU(2)_L\times U(1)_Y,
\label{eq:standard-model-group}
\end{equation}
严格的全局群还可能除以同时作用平凡的离散中心。用 $(R_3,R_2)_Y$ 表示色、弱同位旋表示和超荷，一代左手 Weyl 场可列为
\[
Q_L:(\mathbf3,\mathbf2)_{1/6},\quad
u_R^c:(\bar{\mathbf3},\mathbf1)_{-2/3},\quad
d_R^c:(\bar{\mathbf3},\mathbf1)_{1/3},
\]
\[
L_L:(\mathbf1,\mathbf2)_{-1/2},\qquad
e_R^c:(\mathbf1,\mathbf1)_1 .
\]
这里统一把右手粒子写成左手反粒子，便于反常求和。对表示 $(R_3,R_2)_Y$ 中的场，协变导数为
\[
D_\mu=\partial_\mu-\ii g_s G_\mu^a T_3^a-\ii g W_\mu^i T_2^i-\ii g'YB_\mu .
\]
三个生成元分别只作用在色、弱同位旋和超荷因子上；这正是直积群表示的张量积结构。电荷生成元为（式 \eqnrefs{eq:gell-mann-nishijima}）
\begin{equation}
Q=T_3+Y.
\label{eq:gell-mann-nishijima}
\end{equation}
例如 $Q_L$ 的两个分量分别有 $Q=2/3,-1/3$。

Higgs 场 $H:(\mathbf1,\mathbf2)_{1/2}$ 的势取
\[
V(H)=-\mu^2H^\dagger H+\lambda(H^\dagger H)^2,
\qquad \lambda>0.
\]
当 $\mu^2>0$ 时，真空可选为 $\langle H\rangle=(0,v/\sqrt2)^{\mathsf T}$，从而
\[
SU(2)_L\times U(1)_Y\longrightarrow U(1)_{\mathrm{em}},
\qquad Q\langle H\rangle=0.
\]
一般地，若连续群 $G$ 破缺为稳定子群 $H$，在无规范场且无额外约束时
\begin{equation}
N_{\mathrm{NG}}=\dim G-\dim H.
\label{eq:goldstone-counting}
\end{equation}
按式 \eqnrefs{eq:goldstone-counting}，破缺前 $G=SU(2)_L\times U(1)_Y$ 的生成元数为 $\dim SU(2)_L+\dim U(1)_Y=3+1=4$，真空只保留下 $H=U(1)_{\mathrm{em}}$ 的一个生成元，$\dim H=1$，故 $N_{\mathrm{NG}}=4-1=3$。这三个方向分别成为 $W^\pm$ 与 $Z$ 的纵向极化。规范理论中这些方向被相应规范玻色子吸收；这不是 Goldstone 模“消失”，而是自由度从标量场重新组织到有质量矢量场。

稳定子可以直接算出。$SU(2)$ 基本表示中
$T_3=\operatorname{diag}(1/2,-1/2)$，而 Higgs 超荷为 $Y=1/2$，
所以
\[
(T_3+Y)\langle H\rangle=0.
\]
其余两个 $SU(2)$ 生成元把下分量搬到上分量，
$(T_3-Y)\langle H\rangle\neq0$。因此唯一保持真空不变的连续生成元
是 $Q=T_3+Y$。

把真空代入动能项 $(D_\mu H)^\dagger D^\mu H$，并定义
\[
W_\mu^\pm=\frac{W_\mu^1\mp\ii W_\mu^2}{\sqrt2},
\qquad
\begin{pmatrix}Z_\mu\\A_\mu\end{pmatrix}
=\begin{pmatrix}\cos\theta_W&-\sin\theta_W\\
\sin\theta_W&\cos\theta_W\end{pmatrix}
\begin{pmatrix}W_\mu^3\\B_\mu\end{pmatrix},
\qquad \tan\theta_W=\frac{g'}g,
\]
得到
\[
m_W=\frac{gv}{2},\qquad
m_Z=\frac{v}{2}\sqrt{g^2+g'^2},\qquad
m_A=0.
\]
零质量组合 $A_\mu$ 正对应未破缺生成元 $Q$；三个有质量矢量场
分别吸收三个破缺方向。

表示张量积是否含平凡表示，可以直接检验相互作用是否允许。例如下型 Yukawa 项可写为
\[
\mathcal L_d=-y_d\,\overline Q_{L\,i}^{\ a}H^i d_{R\,a}+\mathrm{h.c.}
\]
色指标 $a$ 由 $\bar{\mathbf3}\otimes\mathbf3$ 缩并，弱指标 $i$ 由 $\bar{\mathbf2}\otimes\mathbf2$ 缩并，而超荷满足 $-1/6+1/2-1/3=0$；三项同时为单态，耦合才规范不变。上型夸克项使用 $\widetilde H=\ii\sigma_2H^*$，其超荷为 $-1/2$。

手征规范理论只有在所有局域规范反常抵消时才一致。三角形图的群因子可写成
\[
\mathcal M^{abc}_{\triangle}\propto
\operatorname{tr}_R\!\left(T^a\{T^b,T^c\}\right),
\]
它使 Ward 恒等式出现量子修正，示意为
\[
(D_\mu J^\mu)^a
=\frac{g^2}{16\pi^2}\,\mathcal A^{abc}
F^b_{\mu\nu}\widetilde F^{c\mu\nu},
\]
其中总反常系数是所有左手 Weyl 场之和
\[
\mathcal A^{abc}=\sum_{f\,\mathrm{left}}
\operatorname{tr}_{R_f}\!\left(T^a\{T^b,T^c\}\right).
\]
以 $[SU(3)]^2U(1)$ 三角形为例，其系数正比于左手场上 $\operatorname{tr}_R(T^a\{T^b,T^c\})$ 对色空间求迹后乘以超荷 $Y$。对 $SU(3)$ 基本表示取 Gell--Mann 归一化 $\operatorname{tr}(T^aT^b)=\tfrac12\delta^{ab}$，反基本表示同号。逐一代入一代左手场：
\[
\underbrace{2\cdot\tfrac16\cdot\tfrac12}_{Q_L:\,2\text{ 个弱分量}}
+\underbrace{(-\tfrac23)\cdot\tfrac12}_{u_R^c}
+\underbrace{\tfrac13\cdot\tfrac12}_{d_R^c}
=\tfrac16-\tfrac13+\tfrac16=0.
\]
这里色迹已经包含在基本表示的 Dynkin 指数 $T(\mathbf3)=1/2$ 中，不能再额外乘以色数；反基本表示具有相同的二次指标。$[SU(3)]^3$、$[SU(2)]^2U(1)$、$[U(1)]^3$ 和引力--$U(1)$ 反常由同样的逐代求和逐项抵消。抵消是整套表示的性质，不能只检查某一个多重态。

其余求和也可以在一页内算完。纯色反常中，基本表示与反基本表示
符号相反；$Q_L$ 有两个弱分量，因此
\[
[SU(3)]^3:\qquad 2-1-1=0.
\]
$SU(2)$ 基本表示的二次指标为 $1/2$。夸克双重态有三种颜色，
故
\[
[SU(2)]^2U(1):\qquad
3\left(\frac16\right)\frac12
+\left(-\frac12\right)\frac12=0.
\]
对纯超荷反常，每个独立 Weyl 分量都要计数：
\[
\begin{aligned}
[U(1)]^3:\quad&
6\left(\frac16\right)^3
+3\left(-\frac23\right)^3
+3\left(\frac13\right)^3\\
&\quad
+2\left(-\frac12\right)^3+1^3=0.
\end{aligned}
\]
引力--$U(1)$ 混合反常把三次方换成一次方，
\[
6\left(\frac16\right)
+3\left(-\frac23\right)
+3\left(\frac13\right)
+2\left(-\frac12\right)+1=0.
\]
$SU(2)$ 的基本表示是伪实表示，所以局域三角反常自动为零；
还需检查全局 Witten 反常。一代中共有三个带颜色标记的夸克
双重态和一个轻子双重态，总数为 $4$，是偶数，故全局反常也消失。
这些等式若缺少任何一个标准模型多重态，通常就不再同时成立。

表示给出生成元在每个场上的作用，张量积中的单态决定可写的
相互作用，Higgs 真空的稳定子决定未破缺电荷，反常迹条件则检查
整套手征表示能否成为一致的量子规范理论。
